% ------------------------------------------------------------
% Front-matter reader's note, shared by all variants. Input
% after \tableofcontents, BEFORE the first \section. Deliberately
% free of \cite commands: citations outside numbered sections
% would corrupt the section-level backrefs (both works mentioned
% here are cited properly in \cref{sec:roundshape}).
% ------------------------------------------------------------
\noindent\textit{A word on how to read what follows.} This is a
deliberately leisurely introduction: we approach \SHA{} the long way
around, with excursions into history, epistemology, finite geometry, and
analytic combinatorics, because the object of interest is not the
algorithm so much as where it sits---its context, its significance, and
the peculiar shape of what is and is not known about it. Readers
impatient for the algorithm itself should know that it is short and
entirely public: the Wikipedia page on SHA-2 contains a complete and
perfectly accessible, if dense, definition in pseudocode, and the
official standard (NIST's FIPS 180-4) is only a few pages longer; both
are cited where the round function is anatomized (\cref{sec:roundshape}).
Nothing about the definition is hidden. Everything about its behavior is.
\bigskip
